% Modelagem do filtro no domínio da frequência
\begin{frame}{Modelagem no domínio da frequência}
  \small
  \begin{columns}[T,onlytextwidth]
    \begin{column}{0.47\textwidth}
      \begin{block}{Equações nodais}
        \[
          V_1=-\frac{V_o}{CR_3s}
        \]
        \[
          \frac{-V_1+V_i}{R_1}
          =CsV_1+Cs(V_1-V_o)+\frac{V_1}{R_2}.
        \]
        Considera-se o amplificador operacional ideal: corrente de entrada nula e
        curto-circuito virtual entre os terminais.
      \end{block}
    \end{column}
    \begin{column}{0.49\textwidth}
      \begin{alertblock}{Função de transferência}
        \[
          H(s)=\frac{V_o(s)}{V_i(s)}
          =\frac{-\dfrac{s}{R_1C}}
          {s^2+\dfrac{2s}{R_3C}
          +\dfrac{R_1+R_2}{R_1R_2R_3C^2}}.
        \]
        O zero na origem e o par de polos complexos produzem a resposta passa-faixa.
      \end{alertblock}
    \end{column}
  \end{columns}

  \vspace{0.12cm}
  \centering
  \alert{Para o regime permanente senoidal, substitui-se $s=j\omega$, com $\omega=2\pi f$.}
\end{frame}